\documentclass[aspectratio=169, 14pt]{beamer}
\usepackage{beamer_preamble}

\begin{document}

% ==== Title ===============================================================
\titleslide{Lesson 5-1 \textperiodcentered\ Single-Period Quick}{nth Roots, Rational Exponents + Cylinder DOK-3 (Item 1A Prep)}

% ==== Do Now ==============================================================
\begin{frame}{Do Now --- \(y = x^2\) graph: roots from squares}{Algebra 2 \textperiodcentered\ Lesson 5-1}
\phasebadges{1}{5}{bridge to nth roots}
\vspace{0.2cm}
\begin{projcard}{Do Now --- Explore \& Reason: \(y = x^2\)}{slidegold}
\small
Graph of \(y = x^2\) is projected. For each, find the missing value:

\vspace{0.15cm}
\begin{tabular}{ll}
(a) \((2,\ \_\_)\) \quad (b) \((3,\ \_\_)\) \quad (c) \((-3,\ \_\_)\) & \quad (d) \((5,\ \_\_)\) \\[4pt]
(e) \((\_\_,\ 4)\) \quad (f) \((\_\_,\ -16)\) & \quad (g) \((\_\_,\ 7)\)
\end{tabular}

\vspace{0.2cm}
\textbf{B.} How do you find \(\sqrt{13}\) using this graph?

\vspace{0.1cm}
\textbf{C.} Sketch \(y = x^3\). Approximate \(\sqrt[3]{5}\), \(\sqrt[3]{-5}\), \(\sqrt[3]{8}\).
\end{projcard}
\end{frame}

% ==== Launch ==============================================================
\begin{frame}{Launch --- Example 1 + Example 3 (two examples, compressed)}{Algebra 2 \textperiodcentered\ Lesson 5-1}
\phasebadges{1--2}{13}{teacher models}
\vspace{0.1cm}
\begin{projcard}{nth Roots \(\leftrightarrow\) Rational Exponents --- Rules Card}{slidegreen}
\small
\begin{enumerate}
  \setlength{\itemsep}{1pt}
  \item \(x^{1/n} = \sqrt[n]{x}\) \quad (principal nth root)
  \item \(x^{m/n} = \bigl(\sqrt[n]{x}\bigr)^m = \sqrt[n]{x^m}\)
  \item Rationalize nth-root denominator: multiply by \(\sqrt[n]{x^{n-1}}\)
  \item Even index: radicand \(\ge 0\) in Reals; odd index: any real
\end{enumerate}

\vspace{0.1cm}
\textcolor{slideaccent}{\textbf{Example 1:}} real nth roots \quad\textbf{|}\quad \textcolor{slideaccent}{\textbf{Example 3:}} evaluate rational exponents\\[2pt]
\small Release to Explore at minute 18. Copy Rules card into margin.
\end{projcard}
\end{frame}

% ==== Explore (items 1-4) =================================================
\begin{frame}[shrink=12]{Explore --- TPS items 1--4 (first 14 min)}{Algebra 2 \textperiodcentered\ Lesson 5-1}
\phasebadges{2--3}{32}{student work}
{\small\textcolor{slidegray}{7 items in order. Plan \(\sim\)10 min for Practice \#50.}}

\vspace{0.1cm}
\begin{projcard}{Try It 1 \textperiodcentered\ Real nth roots}{slideblue}
\small Find the real fourth roots of 81; the real cube roots of 64.
\end{projcard}
\vspace{-0.15cm}
\begin{projcard}{Try It 3 \textperiodcentered\ Evaluate rational exponents}{slideblue}
\small Evaluate \(-(16^{3/4})\) and \(\sqrt[5]{3.5^4}\). Round if needed.
\end{projcard}
\vspace{-0.15cm}
\begin{projcard}{Practice \#28 \textperiodcentered\ nth roots with variables}{slideblue}
\small Find the real square roots of 25.
\end{projcard}
\vspace{-0.15cm}
\begin{projcard}{Practice \#30 \textperiodcentered\ Rational exponent form}{slideblue}
\small Rewrite \(\sqrt[6]{729}\) using a fractional exponent.
\end{projcard}
\end{frame}

% ==== Explore (items 5-7) =================================================
\begin{frame}[shrink=8]{Explore --- TPS items 5--7 + DOK-3 Capstone (18 min)}{Algebra 2 \textperiodcentered\ Lesson 5-1}
\phasebadges{2--3}{32}{student work}
\vspace{0.1cm}
\begin{projcard}{Practice \#37 \textperiodcentered\ Simplify: mixed radical/rational}{slideblue}
\small Simplify \(\sqrt[6]{729a^{24}b^{18}}\). Assume all variables positive.
\end{projcard}
\vspace{-0.15cm}
\begin{projcard}{Practice \#48 \textperiodcentered\ Multi-step rational exponent}{slideblue}
\small Which expressions equal \(b^{3/4}\)? (Yes/No table --- see packet.)
\end{projcard}
\vspace{-0.15cm}
\begin{projcard}{Practice \#50 \textperiodcentered\ CYLINDER MODELING \(\bigstar\)}{slideaccent}
\small Cylindrical container: height \(=\) diameter. Volume \(= 169.65\) ft\(^3\).\\
\textbf{Part A:} lateral surface area? \quad \textbf{Part B:} containers in 20-ft hold?\\
\textit{Set up \(V = \pi r^2(2r)\), isolate \(r\) in rational-exponent form, then rationalize.}
\end{projcard}
\end{frame}

% ==== Share / Summary =====================================================
\begin{frame}[shrink=8]{Share / Summary --- Cylinder Reveal}{Algebra 2 \textperiodcentered\ Lesson 5-1}
\phasebadges{2}{5}{DOK-3 reveal + assessment bridge}
\vspace{0.2cm}
\begin{projcard}{Sentence frame \textperiodcentered\ Cylinder Capstone}{slideblue}
\small
``I isolated \(r\) from \(V = \pi r^2 h\) by substituting \(h = 2r\), giving
\(V =\) \underline{\hspace{1.6cm}}.
Solving for \(r\), I got \(r =\) \underline{\hspace{1.6cm}}.
In rational-exponent form, \(r = \bigl(V/(2\pi)\bigr)^{1/3}\).
To rationalize I multiplied by \underline{\hspace{1cm}} to get \underline{\hspace{1.6cm}}.''
\end{projcard}

\vspace{0.15cm}
\begin{projcard}{Bridge to Topic 5 Assessment --- Item 1A Prep}{slideaccent}
\small
The cylinder move: isolate a variable \(\rightarrow\) rational-exponent form \(\rightarrow\) rationalize.
This is \textbf{exactly} what Topic 5 Assessment Item 1A asks. You have now rehearsed it.
\end{projcard}
\end{frame}

% ==== Exit Ticket =========================================================
\begin{frame}{Exit Ticket --- Summary Recap}{Algebra 2 \textperiodcentered\ Lesson 5-1}
\phasebadges{1--2}{3}{not graded}
\vspace{0.2cm}
\begin{projcard}{Complete on your exit ticket}{slidegold}
\begin{enumerate}
  \setlength{\itemsep}{0.25cm}
  \item \(x^{m/n}\) is the same as \underline{\hspace{5cm}} in radical form.
  \item To rationalize a cube-root denominator I \underline{\hspace{4.2cm}}.
  \item One biggest thing I learned today: \underline{\hspace{4cm}}.
\end{enumerate}
\end{projcard}
\end{frame}

\end{document}
